\minitopic{研读实验室：把AI输出放回认识链}\index{检索增强生成}\index{内容溯源}\index{人机监督}\index{机器知识} “AI说$p$”不是一种单一证据。系统可能检索到原文、抽取数据库字段、分类已有输入、预测数值、生成解释，或根据用户要求续写文本。不同任务的真值条件和错误机制不同。把所有输出都叫“答案”，会让核验方法失去对象。 第一步是识别认识角色。检索结果是候选来源位置，摘要是对来源的压缩主张，分类是模型对输入的标签，预测是条件分布，建议还加入目标与价值。一个系统同时完成多项任务时，应逐层审查；检索正确不保证摘要忠实，摘要忠实也不保证建议合理。 第二步区分系统性能与个案证据。某模型在基准集准确率90\%，不能推出这次输出为真的概率是90\%。基准人群、任务难度、分布变化和置信校准都会影响个案。反过来，一次错误也不证明工具在所有任务无用。理性依赖需要错误类型和适用范围，而非总体崇拜或排斥。

\minitopic{生成、引用与可追溯性}生成模型产生与语境相容的符号序列，不内置“每句话有一个真实出处”的保证\parencite{bender2021}。即使输出包含作者、年份和网址，也可能组合出不存在的文献，或用真实文献支持无关主张。形式完整性会加强流畅性错觉，因此引用越具体，越需要逐项核对。 核对引用至少有四层：作品是否存在；元数据是否准确；来源是否真的包含相关主张；主张强度与来源限定是否一致。只通过第一层不能称引证有效。综述文章适合定位，关键定义、数据和原始论证应回到原文；无法取得时应标明二手转引。 检索增强生成把文档片段提供给模型，可减少某些无依据输出，却不能保证片段选择、理解和综合正确。来源库若过期，答案会稳定地过期；切片丢失上下文，模型可能把例外当一般规则；多个来源冲突时，流畅综合可能隐藏分歧。检索是新证据通道，也引入选择层。 内容凭证可记录文件的创建、编辑和签名历史。C2PA规范把可验证来源声明组织为可检查的内容凭证\parencite{c2pa2025}。它能回答某文件是否与特定签名和编辑链一致，却不能单独证明画面所描绘事件为真，也不能保证缺少凭证的内容为假。来源完整性与内容真实性必须分开。

\minitopic{评估：从平均分到错误地图}系统评估应先定义任务单位和损失。事实问答可检查支持关系，摘要可评遗漏与歪曲，医疗分类需分别计算漏诊与误报，开放式解释还要评价可理解性与适用边界。把异质任务压成一个总分，会让高分没有可操作含义。 数据划分必须防止泄漏。若测试题或近似答案出现在训练数据中，表现测到记忆而非泛化；若同一患者的多次记录跨训练和测试集，模型可能利用身份线索。时间外测试、机构外测试和分群评估分别检验不同迁移。 平均准确率会掩盖分布性失败。两个系统总正确率均90\%，一个对各群体相近，另一个对多数群体98\%、少数群体50\%；部署后果显然不同。公平评价还需结合基础率、错误成本与资源分配，不能只要求所有指标数值相同\parencite{buolamwini2018}。 校准检验模型报告的置信是否与长期频率匹配。系统若对报80\%的事件仅五成正确，就会诱导过度依赖；但校准良好也不保证分辨力、因果理解或个案知识。人机界面应说明置信对象，是下一个词、分类标签、模型间一致性还是事实正确概率。 红队与压力测试主动寻找失败，不应只追求展示最好平均值。对抗输入、领域外样本、语言变体、缺失数据和恶意提示可揭示边界。测试结果要保存版本，因为模型、提示、检索库和安全策略变化后，旧结论未必适用。

\minitopic{人机监督的结构条件}“人工在环”只有在监督者拥有相关信息、时间、能力和否决权时才有认识价值。让操作员每小时批准数百个高复杂度结论，实际上把责任签名化；若界面只显示模型结论而隐藏证据，人类很难形成独立判断。 有效监督可采用分歧优先：把模型低置信、多个模型冲突、规则异常和高损失个案交给人；常规个案抽样审计，以防系统把错误集中在未触发区域。监督者先独立判断再看建议，可减少锚定，但会增加成本。程序应与风险相称。 覆写率不是监督质量的充分指标。人很少覆写，可能模型优秀，也可能人无力反对；经常覆写，可能人纠错，也可能算法厌恶。要按最终结果和理由类型评估：哪些覆写改善表现，哪些重复人为偏见，模型与人是否在同类案例共同失败。 责任必须沿链条分配。开发者负责训练与测试说明，部署机构负责场景适配和监控，专业使用者负责个案判断，管理者负责资源和升级流程。最终签字者不能吸收所有上游责任，系统也不能成为无主体的替罪者。内部算法审计可把这些阶段的决定与记录连接起来\parencite{raji2020}；NIST的生成式AI风险管理资料也强调治理、内容来源、测量和持续管理等跨周期任务\parencite{nist2024}。

\minitopic{认识代理与机器知识}机器能否知道取决于“知道”的理论和归属层级。若知识只要求可靠、对事实敏感的真表征，某些系统可能满足；若要求信念、理解理由、认识责任或第一人称承诺，当前系统是否具备便很有争议。不能从“它能答对”直接推出“它知道”，也不能因非人就先验排除。 工具观点把系统看作人类知识形成的组成部分，知识主要归属于胜任使用者或组织。代理观点认为，复杂系统有内部表征、目标和更新，可成为有限认识主体。社会角色观点则关注我们是否应允许系统断言、作证和承担责任。三种问题不同：形而上状态、认识功能和制度许可。 拟人化语言具有实践后果。说“系统相信你有欺诈风险”可能隐藏是机构选择了阈值；说“AI决定拒贷”可能让设计与政策责任消失。技术描述应尽量写清输入、模型输出、阈值和授权行为。是否授予主体词汇应由解释增益决定，而非界面语气。 即使不承认机器知识，人机系统也可能扩展人类知识；即使承认机器具有某种知识，也不自动获得道德权利或法律责任能力。认识地位、道德地位和法律人格没有简单蕴涵。

\minitopic{数字保存、版本与遗忘}数字信息看似永久，实际依赖格式、服务器、权限和维护。链接腐烂、数据库覆盖和平台关闭会切断证据链。研究工作应保存稳定标识、访问日期、版本、必要的本地副本和许可范围。截图保留表面，却可能丢失交互内容和元数据。 版本控制是一种认识工具。它记录谁在何时改变何处，使读者区分原始主张、更正和后来解释。只保留最终文本会让合理更新看似从未发生，也无法分析错误如何产生。对模型输出，还要记录模型版本、系统提示、检索材料和关键参数，才能有限复现。 “被遗忘权”与档案完整性可能冲突。删除个人数据保护隐私，却可能影响研究复核；永久保存公共错误又会持续伤害。制度需按用途、风险和权利设定访问、匿名化和保留期限，而不是把更多数据视为无条件认识增益。

\minitopic{综合案例：AI辅助法律研究}律师询问系统某司法辖区是否承认一项抗辩。系统给出肯定结论和三份判例。第一份存在但来自另一辖区，第二份已被后案限制，第三份不存在。若律师只检查链接是否能打开，仍会把错误权威链写入意见。 合格流程先把问题拆成法域、时间、程序阶段和事实要素；在权威数据库核对每份判例及后续处理；阅读裁判中与抗辩有关的具体段落；查找反向案例；记录检索日期和未决分歧。生成文本可帮助列关键词，不能替代权威性和有效性核验。 若系统遗漏一份不利先例，原因可能是检索库缺失、排序偏差或提示只要求支持材料。律师有对抗性搜索义务，不能把“模型未返回”当作不存在。机构应抽查遗漏率和法域错误，而非只评价答案语气。 最终文书责任仍由律师和机构承担。若工具供应商隐瞒数据库范围，也负有上游责任；若法院要求披露AI使用，应说明与裁判可靠性有关的部分，不必暴露受保护客户信息。认识透明与其他法律义务需同时满足。

\begin{tcolorbox}[breakable,colback=white,colframe=blue!30!black,title={本章研读任务},fonttitle=\bfseries]
选择一次AI辅助任务，逐层标出检索、抽取、生成、判断和行动；为每层定义真值条件、错误类型和责任人；建立至少十例的小型测试集，报告分组错误与校准；最后设计一套包含版本、来源、覆写理由和复核结果的审计记录。
\end{tcolorbox}
